% FILE: Formal causal hierarchy analysis — trigger sufficiency, lock removal, sensitivity analysis, minimum intervention sets
\chapter{Formal Causal Hierarchy Analysis}
\label{ch:causal-hierarchy-formal}

\epigraph{The purpose of models is not to fit the data but to sharpen the questions.}{Samuel Karlin}

Chapter~\ref{ch:causal-hierarchy} proposed a qualitative three-tier classification of ME/CFS mechanisms: trigger-capable root causes, amplifiers, and load-bearing locks. That classification, though clinically useful, rested on verbal reasoning about feedback topology. Verbal arguments cannot determine whether a proposed mechanism is \emph{sufficient} to trigger disease onset, quantify how much a given lock contributes to disease maintenance, or identify the \emph{minimal} set of interventions required to escape a disease attractor. These are mathematical questions, and they demand mathematical answers.

This chapter provides those answers. Using the 64-variable integrated ODE model developed in Chapter~\ref{ch:integrated-systems}---extended here with three new modules---we formalize the causal hierarchy through three complementary analyses:

\begin{enumerate}
    \item \textbf{Trigger sufficiency analysis}: Can perturbing a single mechanism, starting from the healthy steady state, push the system across the separatrix into a disease attractor?
    \item \textbf{Parameter sensitivity analysis}: Which parameters most strongly control the disease steady state?
    \item \textbf{Lock removal analysis}: Can restoring a single mechanism to healthy values, starting from the disease steady state, allow the system to escape the disease attractor?
\end{enumerate}

Together, these analyses transform the qualitative hierarchy of Chapter~\ref{ch:causal-hierarchy} into a quantitative framework with testable predictions.

\begin{limitation}[Analytical vs.\ Numerical Results]
\label{lim:analytical-vs-numerical}
The results in this chapter derive from analytical properties of the ODE system (steady-state equations, Jacobian structure, implicit function theorem) and from reduced subsystem models amenable to closed-form analysis. Numerical simulation of the full 67-variable system (Section~\ref{sec:new-model-components}) has not been performed. Threshold values for trigger sufficiency and lock removal are estimates from reduced models and should be treated as order-of-magnitude predictions. The \emph{qualitative} predictions---which mechanisms are trigger-capable, which are necessary locks, and which intervention combinations are required---are expected to be robust to model details, because they depend on feedback topology rather than precise parameter values.
\end{limitation}

%% ========================================================================
\section{Methodology}
\label{sec:causal-hierarchy-methodology}
%% ========================================================================

\subsection{Notation and Definitions}

Let $\mathbf{x} \in \mathbb{R}^{n}$ denote the state vector ($n = 67$ after the extensions of Section~\ref{sec:new-model-components}) and $\boldsymbol{\theta} \in \mathbb{R}^{m}$ the parameter vector. The system dynamics are:

\begin{equation}
\label{eq:ode-system-ch33}
\frac{d\mathbf{x}}{dt} = \mathbf{f}(\mathbf{x}, \boldsymbol{\theta})
\end{equation}

We denote:

\begin{itemize}
    \item $\boldsymbol{\theta}_h$: parameter vector corresponding to healthy physiology
    \item $\boldsymbol{\theta}_d$: parameter vector corresponding to established ME/CFS (specific values depend on subtype; see Section~\ref{sec:bifurcation-analysis})
    \item $\mathbf{x}_h^* = \mathbf{x}^*(\boldsymbol{\theta}_h)$: healthy steady state, satisfying $\mathbf{f}(\mathbf{x}_h^*, \boldsymbol{\theta}_h) = \mathbf{0}$
    \item $\mathbf{x}_d^* = \mathbf{x}^*(\boldsymbol{\theta}_d)$: disease steady state (one of the four attractors identified in Section~\ref{sec:bifurcation-analysis})
    \item $\mathcal{B}_h$: basin of attraction of the healthy state under parameters $\boldsymbol{\theta}_h$
    \item $\mathcal{B}_d$: basin of attraction of the disease state under parameters $\boldsymbol{\theta}_d$
    \item $\Sigma_{\mathrm{sep}}$: the separatrix hypersurface separating $\mathcal{B}_h$ from $\mathcal{B}_d$
\end{itemize}

The separatrix $\Sigma_{\mathrm{sep}}$ is the stable manifold of the saddle-type steady state that lies between the healthy and disease attractors. Its codimension is 1 in the state space, meaning it is an $(n-1)$-dimensional surface that partitions $\mathbb{R}^{n}$ into the two basins.

\subsection{Trigger Sufficiency Analysis}
\label{sec:trigger-sufficiency-method}

A mechanism $j$ is \emph{trigger-sufficient} if perturbing its associated parameter $\theta_j$ alone, starting from the healthy steady state, causes the system trajectory to cross $\Sigma_{\mathrm{sep}}$.

\begin{definition}[Trigger Sufficiency]
\label{def:trigger-sufficient}
Parameter $\theta_j$ is trigger-sufficient if there exists a value $\theta_j^{\mathrm{crit}}$ such that, with all other parameters held at healthy values,
\begin{equation}
\label{eq:trigger-sufficiency}
\mathbf{x}_h^*(\boldsymbol{\theta}_h|_{\theta_j \to \theta_j^{\mathrm{crit}}}) \notin \mathcal{B}_h
\end{equation}
That is, the steady state under the perturbed parameter lies outside the basin of attraction of the healthy state. The \emph{trigger threshold} is the minimal perturbation:
\begin{equation}
\label{eq:trigger-threshold}
\Delta\theta_j^{\mathrm{trig}} = \min_{\theta_j^{\mathrm{crit}}} |\theta_j^{\mathrm{crit}} - \theta_{j,h}| \quad \text{such that } \mathbf{x}_h^*(\boldsymbol{\theta}_h|_{\theta_j \to \theta_j^{\mathrm{crit}}}) \notin \mathcal{B}_h
\end{equation}
\end{definition}

Operationally, trigger sufficiency is assessed by numerical continuation of the healthy steady state as $\theta_j$ is varied, monitoring for saddle-node bifurcations (the healthy attractor disappears) or for the trajectory crossing $\Sigma_{\mathrm{sep}}$ (the healthy attractor persists but the system is attracted to the disease state). The former is a stronger condition (the healthy state ceases to exist); the latter is a basin boundary crossing that depends on the full transient dynamics.

\subsection{Sensitivity Analysis}
\label{sec:sensitivity-method}

The normalized sensitivity of state variable $x_i^*$ to parameter $\theta_j$ at the disease steady state is:

\begin{equation}
\label{eq:normalized-sensitivity}
S_{ij} = \frac{\partial x_i^*}{\partial \theta_j} \cdot \frac{\theta_j}{x_i^*}
\end{equation}

This is computed by implicit differentiation of $\mathbf{f}(\mathbf{x}^*, \boldsymbol{\theta}) = \mathbf{0}$. Differentiating with respect to $\theta_j$:

\begin{equation}
\label{eq:implicit-differentiation}
\mathbf{J} \cdot \frac{\partial \mathbf{x}^*}{\partial \theta_j} + \frac{\partial \mathbf{f}}{\partial \theta_j} = \mathbf{0}
\end{equation}

where $\mathbf{J} = \partial \mathbf{f}/\partial \mathbf{x}|_{\mathbf{x}^*}$ is the Jacobian evaluated at the steady state. Provided $\mathbf{J}$ is nonsingular (guaranteed at a stable steady state, where all eigenvalues have negative real parts):

\begin{equation}
\label{eq:sensitivity-solution}
\frac{\partial \mathbf{x}^*}{\partial \theta_j} = -\mathbf{J}^{-1} \cdot \frac{\partial \mathbf{f}}{\partial \theta_j}
\end{equation}

The sensitivity matrix $\mathbf{S} \in \mathbb{R}^{n \times m}$ is then assembled from normalized entries. High $|S_{ij}|$ indicates that parameter $\theta_j$ is \emph{load-bearing} for state variable $x_i$: small changes in $\theta_j$ produce large changes in $x_i^*$.

For global sensitivity, we employ Sobol indices. The total-effect Sobol index $S_j^T$ quantifies the fraction of variance in a scalar output $y = g(\mathbf{x}^*(\boldsymbol{\theta}))$ attributable to parameter $\theta_j$ and all its interactions:

\begin{equation}
\label{eq:sobol-total}
S_j^T = 1 - \frac{\mathrm{Var}_{\boldsymbol{\theta}_{\sim j}}[\mathbb{E}_{\theta_j}[y \mid \boldsymbol{\theta}_{\sim j}]]}{\mathrm{Var}[y]}
\end{equation}

where $\boldsymbol{\theta}_{\sim j}$ denotes all parameters except $\theta_j$. High $S_j^T$ means that $\theta_j$ is influential even accounting for nonlinear interactions.

\subsection{Lock Removal Analysis}
\label{sec:lock-removal-method}

A mechanism $j$ is a \emph{necessary lock} if restoring its parameter to the healthy value, starting from the disease steady state, allows the system to escape $\mathcal{B}_d$.

\begin{definition}[Lock Necessity]
\label{def:lock-necessity}
Parameter $\theta_j$ is a necessary lock for disease attractor $\mathbf{x}_d^*$ if:
\begin{equation}
\label{eq:lock-necessity}
\mathbf{x}_d^*(\boldsymbol{\theta}_d|_{\theta_j \to \theta_{j,h}}) \notin \mathcal{B}_d
\end{equation}
That is, restoring $\theta_j$ alone to its healthy value causes the disease steady state to disappear or to leave the basin of attraction.
\end{definition}

Lock removal is assessed by the reverse of the trigger sufficiency protocol: starting from $\boldsymbol{\theta}_d$, restore $\theta_j$ to $\theta_{j,h}$ while holding all other disease parameters fixed, and check whether the disease attractor persists.

A mechanism that is neither trigger-sufficient nor a necessary lock is classified as an \emph{amplifier}: it worsens the disease state when active and improves it when removed, but it cannot independently cause or maintain the disease.

%% ========================================================================
\section{New Model Components}
\label{sec:new-model-components}
%% ========================================================================

The 64-variable model of Section~\ref{sec:updated-state-vector} captures the core physiology but omits three mechanisms that Chapter~\ref{ch:causal-hierarchy} identifies as candidate root causes: GPCR autoantibodies, TRPM3/calcium signaling dysfunction, and the metabolic safe mode hypothesis. This section extends the model with dedicated modules for each.

\subsection{GPCR Autoantibody Module}
\label{sec:gpcr-model}

Anti-G-protein-coupled receptor (anti-GPCR) autoantibodies have been reported in ME/CFS patients, with elevated titers against $\beta_2$-adrenergic, muscarinic M3/M4, and angiotensin II type~1 receptors~\cite{Loebel2016,Sotzny2021}. These autoantibodies are hypothesized to cause endothelial dysfunction and impaired vasoregulation.

\subsubsection{State Variable}

We introduce a single lumped state variable $[\text{AAb}]$ representing the aggregate anti-GPCR autoantibody titer (units: arbitrary normalized, with $[\text{AAb}] = 0$ in health and $[\text{AAb}] = 1$ at maximal pathological levels).

\subsubsection{Autoantibody Dynamics}

Autoantibody production is driven by long-lived plasma cells ($P_{\ell\ell}$), assumed to reside in bone marrow niches (the ``sanctuary hypothesis''; see Limitation~\ref{lim:analytical-vs-numerical} regarding the analytical nature of this analysis), with clearance following standard immunoglobulin turnover:

\begin{equation}
\label{eq:aab-dynamics}
\frac{d[\text{AAb}]}{dt} = \sigma_{\text{AAb}} \cdot \frac{P_{\ell\ell}}{K_P + P_{\ell\ell}} - \delta_{\text{AAb}} \cdot [\text{AAb}]
\end{equation}

where $\sigma_{\text{AAb}}$ is the maximal production rate, $K_P$ is the half-saturation constant for plasma cell--driven production, $P_{\ell\ell}$ is the long-lived plasma cell population (treated as a slowly varying parameter rather than a dynamical variable, reflecting the sanctuary hypothesis that these cells persist independently of ongoing immune activation), and $\delta_{\text{AAb}}$ is the clearance rate (IgG half-life $\approx 21$\,days, so $\delta_{\text{AAb}} \approx 0.033$\,day$^{-1}$).

\subsubsection{Endothelial Dysfunction Coupling}

Autoantibodies impair endothelial function, reducing effective oxygen delivery (Equation~\ref{eq:o2-delivery}):

\begin{equation}
\label{eq:aab-endothelial}
\text{DO}_{2}^{\text{eff}} = \text{DO}_2 \cdot \frac{K_{\text{endo}}^2}{K_{\text{endo}}^2 + [\text{AAb}]^2}
\end{equation}

The Hill coefficient of 2 reflects cooperative binding of autoantibodies to endothelial receptors. At $[\text{AAb}] = K_{\text{endo}}$, oxygen delivery is reduced by 50\%. This effective oxygen delivery replaces $\text{DO}_2$ in the oxygen consumption equation (Equation~\ref{eq:o2-consumption}), propagating the autoantibody effect to mitochondrial ATP production.

\subsubsection{Inflammatory Coupling}

Autoantibody--receptor binding activates monocytes and drives pro-inflammatory cytokine production:

\begin{equation}
\label{eq:aab-inflammation}
\sigma_{\text{cyto,AAb}} = \sigma_{\text{AAb,cyto}} \cdot \frac{[\text{AAb}]}{K_{\text{AAb,cyto}} + [\text{AAb}]}
\end{equation}

This term is added to the cytokine production rates in the immune model (Chapter~\ref{ch:immune-system-models}), providing a tonic inflammatory drive proportional to autoantibody titer. The Michaelis--Menten form ensures saturation at high titers.

\subsubsection{Parameter Table}

\begin{table}[htbp]
\centering
\caption{GPCR autoantibody module parameters.}
\label{tab:aab-parameters}
\begin{tabular}{llll}
\hline
Parameter & Description & Estimated value & Source \\
\hline
$\sigma_{\text{AAb}}$ & Max production rate & 0.05\,day$^{-1}$ & Estimated \\
$K_P$ & Plasma cell half-saturation & 100\,cells/$\mu$L & Estimated \\
$\delta_{\text{AAb}}$ & IgG clearance rate & 0.033\,day$^{-1}$ & IgG half-life \\
$K_{\text{endo}}$ & Endothelial half-inhibition & 0.4 (normalized) & Estimated \\
$\sigma_{\text{AAb,cyto}}$ & Max cytokine induction & 0.1\,pg/mL/day & Estimated \\
$K_{\text{AAb,cyto}}$ & Cytokine induction half-sat & 0.3 (normalized) & Estimated \\
$P_{\ell\ell}$ & Long-lived plasma cells & 50--200\,cells/$\mu$L & ~\cite{Fluge2025daratumumab} \\
\hline
\end{tabular}
\end{table}

\subsection{TRPM3/Calcium Signaling Module}
\label{sec:trpm3-model}

Reduced expression and function of TRPM3 ion channels have been reported in NK cells of ME/CFS patients~\cite{Cabanas2021,Sasso2026trpm3}. TRPM3 mediates calcium entry, and impaired calcium signaling affects both mitochondrial function and immune cell cytotoxicity.

\subsubsection{State Variable}

We introduce $[\text{Ca}^{2+}]_{\text{eff}}$, a lumped effective intracellular calcium concentration representing the calcium available for signaling (distinct from total cytoplasmic calcium, which is buffered). The variable is normalized so that $[\text{Ca}^{2+}]_{\text{eff}} = 1$ in health.

\subsubsection{Calcium Dynamics}

The effective intracellular calcium depends on TRPM3-mediated entry, temperature sensitivity, and standard calcium buffering:

\begin{equation}
\label{eq:calcium-dynamics}
\frac{d[\text{Ca}^{2+}]_{\text{eff}}}{dt} = J_{\text{Ca}} - k_{\text{pump}} \cdot [\text{Ca}^{2+}]_{\text{eff}} + J_{\text{release}} - J_{\text{uptake}}
\end{equation}

where the TRPM3-dependent calcium entry is:

\begin{equation}
\label{eq:trpm3-calcium}
J_{\text{Ca}} = g_{\text{TRPM3}} \cdot (V_m - E_{\text{Ca}}) \cdot f_{\text{temp}}(T)
\end{equation}

Here $g_{\text{TRPM3}}$ is the TRPM3 channel conductance (reduced in ME/CFS), $V_m$ is the membrane potential, $E_{\text{Ca}}$ is the calcium reversal potential, and $f_{\text{temp}}(T)$ is a temperature-dependent gating function:

\begin{equation}
\label{eq:trpm3-temperature}
f_{\text{temp}}(T) = \frac{1}{1 + e^{-\lambda(T - T_{1/2})}}
\end{equation}

with $T_{1/2} \approx 37\,{}^{\circ}\text{C}$ and $\lambda$ determining the steepness of temperature activation. The terms $J_{\text{release}}$ and $J_{\text{uptake}}$ represent ER calcium release and SERCA pump uptake, modeled as:

\begin{align}
\label{eq:calcium-stores}
J_{\text{release}} &= k_{\text{IP3}} \cdot [\text{IP}_3] \cdot ([\text{Ca}^{2+}]_{\text{ER}} - [\text{Ca}^{2+}]_{\text{eff}}) \\
J_{\text{uptake}} &= V_{\text{SERCA}} \cdot \frac{[\text{Ca}^{2+}]_{\text{eff}}^2}{K_{\text{SERCA}}^2 + [\text{Ca}^{2+}]_{\text{eff}}^2}
\end{align}

For the causal hierarchy analysis, we treat the ER store dynamics as fast and use the quasi-steady-state approximation $J_{\text{release}} \approx J_{\text{uptake}}$ (justified by the rapid timescale of ER calcium equilibration, typically seconds, relative to the minutes-to-hours timescale of TRPM3-mediated transmembrane calcium entry), reducing the calcium equation to:

\begin{equation}
\label{eq:calcium-qss}
\frac{d[\text{Ca}^{2+}]_{\text{eff}}}{dt} \approx J_{\text{Ca}} - k_{\text{pump}} \cdot [\text{Ca}^{2+}]_{\text{eff}}
\end{equation}

\subsubsection{Mitochondrial Coupling}

Calcium modulates Complex~I activity through allosteric regulation of matrix dehydrogenases:

\begin{equation}
\label{eq:calcium-mito}
\alpha_{\text{CI}}^{\text{eff}} = \alpha_{\text{CI}} \cdot \left[(1 - \varepsilon_{\text{Ca}}) + \varepsilon_{\text{Ca}} \cdot \frac{[\text{Ca}^{2+}]_{\text{eff}}^2}{K_{\text{Ca}}^2 + [\text{Ca}^{2+}]_{\text{eff}}^2}\right]
\end{equation}

where $\varepsilon_{\text{Ca}} \in [0, 1]$ is the fraction of Complex~I activity that is calcium-dependent. At $[\text{Ca}^{2+}]_{\text{eff}} = 0$, the effective Complex~I activity is reduced to $\alpha_{\text{CI}} \cdot (1 - \varepsilon_{\text{Ca}})$; at saturating calcium, full activity is restored. The Hill coefficient of 2 captures the cooperative calcium dependence of NADH dehydrogenase activation. This replaces $\alpha_{\text{CI}}$ in the Complex~I equation (Equation~\ref{eq:complex-i}).

\subsubsection{Immune Coupling}

NK cell cytotoxicity depends on calcium-dependent degranulation. The effective kill rate is:

\begin{equation}
\label{eq:calcium-nk}
k_{\text{kill}}^{\text{eff}} = k_{\text{kill}} \cdot \frac{[\text{Ca}^{2+}]_{\text{eff}}^2}{K_{\text{Ca,NK}}^2 + [\text{Ca}^{2+}]_{\text{eff}}^2}
\end{equation}

Reduced $[\text{Ca}^{2+}]_{\text{eff}}$ (due to impaired TRPM3) simultaneously impairs NK cell function and reduces mitochondrial ATP production---a dual defect from a single channel dysfunction.

\subsubsection{Parameter Table}

\begin{table}[htbp]
\centering
\caption{TRPM3/calcium signaling module parameters.}
\label{tab:trpm3-parameters}
\begin{tabular}{llll}
\hline
Parameter & Description & Estimated value & Source \\
\hline
$g_{\text{TRPM3}}$ & TRPM3 conductance & 1.0 (healthy), 0.3--0.5 (ME/CFS) & ~\cite{Sasso2026trpm3} \\
$V_m$ & Resting membrane potential & $-70$\,mV & Standard \\
$E_{\text{Ca}}$ & Calcium reversal potential & $+120$\,mV & Standard \\
$T_{1/2}$ & Temperature half-activation & $37\,{}^{\circ}\text{C}$ & Estimated \\
$\lambda$ & Temperature steepness & 1.0\,${}^{\circ}\text{C}^{-1}$ & Estimated \\
$k_{\text{pump}}$ & Calcium extrusion rate & 0.5\,s$^{-1}$ & Standard \\
$\varepsilon_{\text{Ca}}$ & Ca-dependent CI fraction & 0.3 & Estimated \\
$K_{\text{Ca}}$ & Mito calcium half-activation & 0.5 (normalized) & Estimated \\
$K_{\text{Ca,NK}}$ & NK calcium half-activation & 0.4 (normalized) & Estimated \\
\hline
\end{tabular}
\end{table}

\subsection{Metabolic Safe Mode Switch}
\label{sec:safe-mode-model}

The ``metabolic safe mode'' hypothesis proposes that ME/CFS represents a deliberate, evolutionarily conserved suppression of energy production in response to perceived threat---analogous to the dauer state in \emph{C.\ elegans} or torpor in hibernating mammals~\cite{Naviaux2016metabolomics,Hrvatin2020torpor}. Rather than mitochondrial \emph{damage}, this model posits an active \emph{downregulation} that persists because the threat signal is not adequately resolved.

\subsubsection{State Variable}

We introduce $\Sigma \in [0, 1]$, representing the activation state of the metabolic safe mode switch. $\Sigma = 0$ corresponds to normal metabolic operation; $\Sigma = 1$ corresponds to full safe mode activation.

\subsubsection{Bistable Switch Dynamics}

The safe mode switch exhibits hysteresis through a bistable ODE:

\begin{equation}
\label{eq:safe-mode-dynamics}
\frac{d\Sigma}{dt} = k_{\text{on}} \cdot \frac{\mathcal{T}^{n_{\text{on}}}}{K_{\mathcal{T}}^{n_{\text{on}}} + \mathcal{T}^{n_{\text{on}}}} \cdot (1 - \Sigma) - k_{\text{off}} \cdot \frac{(1 - \Sigma)^{n_{\text{off}}}}{K_{\text{off}}^{n_{\text{off}}} + (1 - \Sigma)^{n_{\text{off}}}} \cdot \Sigma
\end{equation}

where $\mathcal{T}$ is the aggregate threat signal:

\begin{equation}
\label{eq:threat-signal}
\mathcal{T} = w_{\text{cyto}} \cdot \boldsymbol{C}_{\text{pro}} + w_{\text{ROS}} \cdot [\text{ROS}] + w_{\text{LPS}} \cdot [\text{LPS}] + w_V \cdot V
\end{equation}

with $\boldsymbol{C}_{\text{pro}} = [\text{IL-1}\beta] + [\text{IL-6}] + [\text{TNF-}\alpha]$ the aggregate pro-inflammatory cytokine level, and $w_{\text{cyto}}, w_{\text{ROS}}, w_{\text{LPS}}, w_V$ weighting coefficients.

The key feature of Equation~\ref{eq:safe-mode-dynamics} is the \emph{cooperative self-inhibition} in the off-rate: the term $(1 - \Sigma)^{n_{\text{off}}}/(K_{\text{off}}^{n_{\text{off}}} + (1 - \Sigma)^{n_{\text{off}}})$ decreases as $\Sigma$ increases, meaning that the deeper the system is in safe mode, the harder it is to exit. With $n_{\text{on}} \geq 2$ and $n_{\text{off}} \geq 2$, this creates a bistable switch with distinct activation and deactivation thresholds for $\mathcal{T}$, producing hysteresis.

\subsubsection{Hysteresis Analysis}

At steady state ($d\Sigma/dt = 0$), the activation and deactivation thresholds can be found analytically in the limit of strong cooperativity ($n_{\text{on}}, n_{\text{off}} \gg 1$):

\begin{align}
\label{eq:safe-mode-thresholds}
\mathcal{T}_{\text{on}} &\approx K_{\mathcal{T}} \cdot \left(\frac{k_{\text{off}}}{k_{\text{on}}}\right)^{1/n_{\text{on}}} \\
\mathcal{T}_{\text{off}} &\approx K_{\mathcal{T}} \cdot \left(\frac{k_{\text{off}} \cdot K_{\text{off}}^{n_{\text{off}}}}{k_{\text{on}}}\right)^{1/n_{\text{on}}}
\end{align}

These approximations follow from setting $d\Sigma/dt = 0$ in Equation~\ref{eq:safe-mode-dynamics} and solving for $\mathcal{T}$ in the strong cooperativity limit where the Hill functions approach step functions. The $\mathcal{T}_{\text{on}}$ expression corresponds to the lower saddle-node bifurcation ($\Sigma$ jumping from $\approx 0$ to $\approx 1$), while $\mathcal{T}_{\text{off}}$ corresponds to the upper saddle-node ($\Sigma$ falling from $\approx 1$ to $\approx 0$). The asymmetry $\mathcal{T}_{\text{off}} < \mathcal{T}_{\text{on}}$ arises because the $K_{\text{off}}^{n_{\text{off}}}$ factor in the deactivation term raises the effective threshold for the off-switch.

The hysteresis width $\Delta\mathcal{T} = \mathcal{T}_{\text{on}} - \mathcal{T}_{\text{off}}$ determines how far the threat signal must drop below the activation threshold to deactivate safe mode. The cooperative off-rate ensures $\mathcal{T}_{\text{off}} < \mathcal{T}_{\text{on}}$: safe mode activates at a higher threat level than the level at which it deactivates, creating a ``sticky'' switch that resists deactivation.

\subsubsection{Effects on Energy Metabolism}

When $\Sigma > 0$, the safe mode switch suppresses electron transport chain (ETC) capacity:

\begin{equation}
\label{eq:safe-mode-etc}
V_{\max,k}^{\text{eff}} = V_{\max,k} \cdot (1 - \sigma_{\text{suppress}} \cdot \Sigma), \qquad k \in \{\text{CI}, \text{CIII}, \text{CIV}\}
\end{equation}

where $\sigma_{\text{suppress}} \in [0, 1]$ is the maximal suppression fraction ($\sigma_{\text{suppress}} \approx 0.5$ for 50\% suppression at full safe mode activation).

Additionally, safe mode diverts Krebs cycle flux toward itaconate production, an immunometabolic switch observed in activated macrophages:

\begin{equation}
\label{eq:safe-mode-itaconate}
J_{\text{Krebs}}^{\text{eff}} = J_{\text{Krebs}} \cdot (1 - \sigma_{\text{itaconate}} \cdot \Sigma)
\end{equation}

where $\sigma_{\text{itaconate}}$ quantifies the fraction of Krebs cycle flux diverted to itaconate at full safe mode activation.

\subsubsection{Parameter Table}

\begin{table}[htbp]
\centering
\caption{Metabolic safe mode module parameters.}
\label{tab:safe-mode-parameters}
\begin{tabular}{llll}
\hline
Parameter & Description & Estimated value & Source \\
\hline
$k_{\text{on}}$ & Safe mode activation rate & 0.1\,day$^{-1}$ & Estimated \\
$k_{\text{off}}$ & Safe mode deactivation rate & 0.02\,day$^{-1}$ & Estimated \\
$K_{\mathcal{T}}$ & Threat half-activation & 0.5 (normalized) & Estimated \\
$n_{\text{on}}$ & Activation cooperativity & 3 & Estimated \\
$n_{\text{off}}$ & Deactivation cooperativity & 3 & Estimated \\
$K_{\text{off}}$ & Deactivation half-inhibition & 0.3 & Estimated \\
$\sigma_{\text{suppress}}$ & Max ETC suppression & 0.5 & Estimated \\
$\sigma_{\text{itaconate}}$ & Max Krebs diversion & 0.3 & Estimated \\
$w_{\text{cyto}}$ & Cytokine threat weight & 0.4 & Estimated \\
$w_{\text{ROS}}$ & ROS threat weight & 0.3 & Estimated \\
$w_{\text{LPS}}$ & LPS threat weight & 0.2 & Estimated \\
$w_V$ & Viral threat weight & 0.1 & Estimated \\
\hline
\end{tabular}
\end{table}

\begin{limitation}[Safe Mode Switch Formulation]
\label{lim:safe-mode-switch-formulation}
The bistable safe mode switch (Equation~\ref{eq:safe-mode-dynamics}) is a mathematical hypothesis with no direct experimental validation of the specific functional form, Hill exponents, or parameter values. The qualitative property of interest---hysteretic bistability driven by inflammatory threat signals---is biologically motivated by the observed persistence of metabolic suppression after infection resolution, but the specific ODE formulation is one of many possible mathematical representations of this behavior.
\end{limitation}

\subsection{Updated State Vector}
\label{sec:extended-state-vector}

The three new modules add three state variables to the 64-variable model of Section~\ref{sec:updated-state-vector}:

\begin{enumerate}
    \item $[\text{AAb}]$: GPCR autoantibody titer (Section~\ref{sec:gpcr-model})
    \item $[\text{Ca}^{2+}]_{\text{eff}}$: effective intracellular calcium (Section~\ref{sec:trpm3-model})
    \item $\Sigma$: metabolic safe mode activation (Section~\ref{sec:safe-mode-model})
\end{enumerate}

The extended system comprises $n = 67$ state variables governed by 67 coupled ODEs. The new variables couple to the existing model through the oxygen delivery equation (Equation~\ref{eq:aab-endothelial}), Complex~I activity (Equation~\ref{eq:calcium-mito}), NK cell kill rate (Equation~\ref{eq:calcium-nk}), ETC capacity (Equation~\ref{eq:safe-mode-etc}), and Krebs cycle flux (Equation~\ref{eq:safe-mode-itaconate}). The coupling is unidirectional for autoantibodies (driven by plasma cells, affects energy and cytokines) but bidirectional for the safe mode switch (driven by cytokines and ROS, which are themselves affected by the energy suppression). The TRPM3/calcium module has primarily feedforward effects (reduced calcium impairs both mitochondria and NK cells) with a weak feedback through the energy--immune coupling loop.

%% ========================================================================
\section{Trigger Sufficiency Analysis}
\label{sec:trigger-sufficiency}
%% ========================================================================

We now apply the trigger sufficiency framework (Section~\ref{sec:trigger-sufficiency-method}) to each candidate root cause identified in Chapter~\ref{ch:causal-hierarchy}. In each case, we start from the healthy steady state $\mathbf{x}_h^*$ and perturb a single parameter, asking whether the perturbation is sufficient to drive the system across the separatrix $\Sigma_{\mathrm{sep}}$.

\subsection{Protocol}
\label{sec:trigger-protocol}

For each candidate mechanism with associated parameter $\theta_j$:

\begin{enumerate}
    \item Set all parameters to healthy values $\boldsymbol{\theta}_h$.
    \item Compute the healthy steady state $\mathbf{x}_h^*$.
    \item Vary $\theta_j$ from $\theta_{j,h}$ toward $\theta_{j,d}$ while holding all other parameters fixed.
    \item At each value of $\theta_j$, compute the perturbed steady state $\mathbf{x}^*(\boldsymbol{\theta}_h|_{\theta_j})$ via numerical continuation.
    \item Classify the outcome:
    \begin{itemize}
        \item \emph{Trigger-capable}: a saddle-node bifurcation destroys the healthy attractor, or the perturbed trajectory is captured by a disease attractor.
        \item \emph{Not trigger-capable}: the healthy attractor persists and remains globally attracting for all physiologically plausible values of $\theta_j$.
    \end{itemize}
    \item Record the critical threshold $\theta_j^{\mathrm{crit}}$, the cascade sequence (which state variables deviate first), and the time to reach the disease attractor.
\end{enumerate}

\subsection{CNS Energy Crisis: Complex~I Reduction}
\label{sec:trigger-ci}

\textbf{Parameter}: $\alpha_{\text{CI}}$ (Complex~I activity), varied from $\alpha_{\text{CI}} = 1.0$ (healthy) to $\alpha_{\text{CI}} = 0.0$ (complete Complex~I failure).

\textbf{Reduced model analysis.} To derive the trigger threshold analytically, we reduce the 67-variable system to its essential energy--immune core: a 2-variable system in $[\text{ATP}]$ and $N_a$ (activated immune cells). All other variables are treated as fast (equilibrating on shorter timescales) or slow (approximately constant on the disease-onset timescale). In this reduced model, the healthy steady state satisfies:

\begin{equation}
\label{eq:atp-ss-ci}
[\text{ATP}]^*(\alpha_{\text{CI}}) = \frac{J_{\text{production}}(\alpha_{\text{CI}}) - J_{\text{immune}}^*}{J_{\text{basal}} + J_{\text{activity}}} \cdot [\text{ATP}]_0
\end{equation}

where $J_{\text{production}}$ depends approximately linearly on $\alpha_{\text{CI}}$ near the healthy operating point: $J_{\text{production}}(\alpha_{\text{CI}}) \approx J_{\text{production,max}} \cdot \alpha_{\text{CI}}$. The immune activation rate $k_{\text{act}}^{\text{eff}}$ (Equation~\ref{eq:immune-energy-feedback}) creates a positive feedback loop: as $\alpha_{\text{CI}}$ decreases, $[\text{ATP}]^*$ falls, which impairs immune control, which increases immune activation, which further drains ATP. The immune energy demand at steady state is:

\begin{equation}
\label{eq:immune-demand-ss}
J_{\text{immune}}^* = e_a \cdot N_a^* + e_M \cdot M_a^* + e_T \cdot T_a^*
\end{equation}

where $N_a^*$, $M_a^*$, and $T_a^*$ are the steady-state activated immune cell counts, themselves functions of $[\text{ATP}]^*$ through the Hill-type coupling (Equation~\ref{eq:immune-energy-feedback}). The coupled steady-state equations form a fixed-point problem whose solutions can be analyzed graphically: the energy nullcline ($d[\text{ATP}]/dt = 0$) and the immune nullcline ($dN_a/dt = 0$) intersect at 1, 2, or 3 fixed points depending on $\alpha_{\text{CI}}$. The saddle-node bifurcation---where the healthy and saddle fixed points collide and annihilate---occurs when the two nullclines become tangent.

From the reduced 2-variable energy--immune model, the critical threshold is:

\begin{equation}
\label{eq:ci-critical}
\alpha_{\text{CI}}^{\mathrm{crit}} \approx 1 - \frac{K_{\text{ATP,immune}}^2}{[\text{ATP}]_0^2} \cdot \frac{J_{\text{basal}}}{J_{\text{production,max}}} \cdot \frac{1}{1 - J_{\text{immune,max}}/J_{\text{production,max}}}
\end{equation}

This expression is derived from the tangency condition at the saddle-node bifurcation of the reduced two-variable energy--immune model (Equations~\ref{eq:energy-immune-coupling}--\ref{eq:immune-energy-feedback} of Chapter~\ref{ch:integrated-systems}), under the assumption that ATP production scales linearly with $\alpha_{\text{CI}}$. The specific algebraic form should be treated as an approximation that may be refined by numerical bifurcation analysis of the full system.

Substituting estimated parameter values yields $\alpha_{\text{CI}}^{\mathrm{crit}} \approx 0.65\text{--}0.70$. Below this threshold, the healthy attractor undergoes a saddle-node bifurcation and disappears: the only remaining attractor is the metabolic-dominant disease state.

\textbf{Cascade sequence.} As $\alpha_{\text{CI}}$ crosses the critical threshold:
\begin{enumerate}
    \item $[\text{ATP}]$ drops below the immune cooperation threshold ($t \sim$ hours)
    \item $k_{\text{act}}^{\text{eff}}$ decreases, viral/pathogen control weakens ($t \sim$ days)
    \item Cytokine levels rise, increasing $J_{\text{immune}}$ ($t \sim$ days to weeks)
    \item The energy--immune vicious cycle locks in ($t \sim$ weeks)
    \item Epigenetic consolidation begins ($t \sim$ months, Equation~\ref{eq:epigenetic-evolution})
\end{enumerate}

\textbf{Verdict}: \textbf{Trigger-sufficient}. Complex~I reduction below $\alpha_{\text{CI}}^{\mathrm{crit}} \approx 0.65$ is sufficient to cause disease onset in the reduced energy--immune model without any other parameter perturbation.

\subsection{Metabolic Safe Mode Activation}
\label{sec:trigger-safe-mode}

\textbf{Parameter}: aggregate threat signal $\mathcal{T}$ (Equation~\ref{eq:threat-signal}), representing any source of danger signaling (infection, toxin exposure, psychological stress).

\textbf{Analysis.} The safe mode switch (Equation~\ref{eq:safe-mode-dynamics}) is designed to exhibit hysteresis. Starting from $\Sigma = 0$ (healthy), increasing $\mathcal{T}$ above the activation threshold $\mathcal{T}_{\text{on}}$ (Equation~\ref{eq:safe-mode-thresholds}) triggers the transition to $\Sigma \approx 1$. Once activated, the cooperative self-inhibition in the off-rate ensures that $\Sigma$ remains elevated even when $\mathcal{T}$ drops below $\mathcal{T}_{\text{on}}$, provided $\mathcal{T}$ remains above $\mathcal{T}_{\text{off}} < \mathcal{T}_{\text{on}}$.

The critical feature is the coupling from $\Sigma$ back to $\mathcal{T}$: ETC suppression (Equation~\ref{eq:safe-mode-etc}) increases ROS production (due to backed-up electron carriers), which contributes to $\mathcal{T}$ through the $w_{\text{ROS}} \cdot [\text{ROS}]$ term. Similarly, reduced ATP impairs immune function, potentially increasing viral load $V$ and cytokine levels $\boldsymbol{C}_{\text{pro}}$, further elevating $\mathcal{T}$. This creates a self-sustaining loop: $\Sigma \uparrow \to \text{ETC} \downarrow \to \text{ROS} \uparrow \to \mathcal{T} \uparrow \to \Sigma$ stays elevated.

\textbf{Phase portrait.} The $(\mathcal{T}, \Sigma)$ phase plane exhibits the classical hysteresis loop. The activation trajectory (increasing $\mathcal{T}$) follows the low-$\Sigma$ branch until $\mathcal{T} = \mathcal{T}_{\text{on}}$, then jumps to the high-$\Sigma$ branch. The deactivation trajectory (decreasing $\mathcal{T}$) follows the high-$\Sigma$ branch until $\mathcal{T} = \mathcal{T}_{\text{off}}$. For ME/CFS, the self-sustaining feedback maintains $\mathcal{T} > \mathcal{T}_{\text{off}}$ even after the external trigger resolves, trapping the system in the high-$\Sigma$ state.

\textbf{Cascade sequence.}
\begin{enumerate}
    \item External threat (e.g., viral infection) elevates $\mathcal{T}$ above $\mathcal{T}_{\text{on}}$ ($t \sim$ days)
    \item $\Sigma$ switches to high state, ETC suppressed ($t \sim$ days)
    \item ROS increases, cytokines rise due to impaired immune control ($t \sim$ days to weeks)
    \item Self-sustaining $\mathcal{T}$ maintains $\Sigma$ after external trigger clears ($t \sim$ weeks)
    \item Energy--immune vicious cycle engaged secondarily ($t \sim$ weeks to months)
\end{enumerate}

\textbf{Verdict}: \textbf{Trigger-sufficient}. A transient threat signal exceeding $\mathcal{T}_{\text{on}}$ is sufficient to activate the safe mode switch, which then self-sustains through endogenous feedback loops. The time course matches the clinical observation that ME/CFS onset often follows an acute infection by days to weeks~\cite{Hickie2006postinfectious}.

\subsection{GPCR Autoantibodies}
\label{sec:trigger-aab}

\textbf{Parameter}: $P_{\ell\ell}$ (long-lived plasma cell population) or equivalently $\sigma_{\text{AAb}}$ (autoantibody production rate).

\textbf{Analysis.} The autoantibody module (Equation~\ref{eq:aab-dynamics}) reaches steady state at:

\begin{equation}
\label{eq:aab-steady-state}
[\text{AAb}]^* = \frac{\sigma_{\text{AAb}}}{\delta_{\text{AAb}}} \cdot \frac{P_{\ell\ell}}{K_P + P_{\ell\ell}}
\end{equation}

This steady-state titer determines the oxygen delivery impairment (Equation~\ref{eq:aab-endothelial}) and inflammatory drive (Equation~\ref{eq:aab-inflammation}). The trigger sufficiency question reduces to: does $[\text{AAb}]^*$ alone---with all other parameters healthy---produce sufficient oxygen delivery impairment and inflammatory drive to push the energy--immune system past the separatrix?

From the reduced energy model, the critical autoantibody level satisfies:

\begin{equation}
\label{eq:aab-critical}
\text{DO}_{2}^{\text{eff}}([\text{AAb}]^{\mathrm{crit}}) = \text{DO}_2 \cdot \frac{K_{\text{endo}}^2}{K_{\text{endo}}^2 + ([\text{AAb}]^{\mathrm{crit}})^2} \leq \text{DO}_{2}^{\mathrm{min}}
\end{equation}

where $\text{DO}_{2}^{\mathrm{min}}$ is the oxygen delivery below which the energy--immune bifurcation occurs. With the estimated parameters, $[\text{AAb}]^{\mathrm{crit}} \approx 0.55$--$0.65$ (normalized), corresponding to moderately elevated autoantibody titers.

\textbf{Cascade sequence.}
\begin{enumerate}
    \item Autoimmune event generates $P_{\ell\ell}$, $[\text{AAb}]$ rises ($t \sim$ weeks to months)
    \item Endothelial dysfunction reduces $\text{DO}_{2}^{\text{eff}}$ ($t \sim$ concurrent)
    \item Reduced oxygen delivery impairs mitochondrial function ($t \sim$ concurrent)
    \item $[\text{ATP}]$ falls, engaging the energy--immune vicious cycle ($t \sim$ weeks)
    \item Inflammatory coupling amplifies: $\sigma_{\text{cyto,AAb}}$ drives cytokine elevation ($t \sim$ weeks)
    \item Epigenetic consolidation ($t \sim$ months)
\end{enumerate}

\textbf{Verdict}: \textbf{Trigger-sufficient}. Autoantibody production above $[\text{AAb}]^{\mathrm{crit}}$ is sufficient to cause disease onset through oxygen delivery impairment and inflammatory drive. The slow timescale of autoantibody accumulation (weeks to months) is consistent with the gradual-onset ME/CFS subtype.

\subsection{TRPM3/Calcium Channel Dysfunction}
\label{sec:trigger-trpm3}

\textbf{Parameter}: $g_{\text{TRPM3}}$ (TRPM3 channel conductance), varied from $g_{\text{TRPM3}} = 1.0$ (healthy) to $g_{\text{TRPM3}} = 0.0$ (complete channel loss).

\textbf{Analysis.} Reduced $g_{\text{TRPM3}}$ decreases $[\text{Ca}^{2+}]_{\text{eff}}$ (Equation~\ref{eq:calcium-qss}), which simultaneously reduces Complex~I activity (Equation~\ref{eq:calcium-mito}) and NK cell function (Equation~\ref{eq:calcium-nk}). The effective Complex~I activity becomes:

\begin{equation}
\label{eq:ci-eff-trpm3}
\alpha_{\text{CI}}^{\text{eff}}(g_{\text{TRPM3}}) = \alpha_{\text{CI}} \cdot \left[(1 - \varepsilon_{\text{Ca}}) + \varepsilon_{\text{Ca}} \cdot \frac{[\text{Ca}^{2+}]_{\text{eff}}^{*2}(g_{\text{TRPM3}})}{K_{\text{Ca}}^2 + [\text{Ca}^{2+}]_{\text{eff}}^{*2}(g_{\text{TRPM3}})}\right]
\end{equation}

For TRPM3 alone to be trigger-sufficient, the reduced $\alpha_{\text{CI}}^{\text{eff}}$ must fall below $\alpha_{\text{CI}}^{\mathrm{crit}} \approx 0.65$ (Section~\ref{sec:trigger-ci}). With $\varepsilon_{\text{Ca}} = 0.3$, even complete loss of TRPM3 ($g_{\text{TRPM3}} = 0$) reduces $\alpha_{\text{CI}}^{\text{eff}}$ to $\alpha_{\text{CI}} \cdot 0.70$---only reaching the bifurcation threshold if baseline $\alpha_{\text{CI}}$ is already near 0.93. In most individuals with normal baseline Complex~I activity ($\alpha_{\text{CI}} = 1.0$), the reduction to $0.70$ is insufficient to cross the saddle-node bifurcation.

However, the simultaneous impairment of NK cell function ($k_{\text{kill}}^{\text{eff}} \to 0$ at $g_{\text{TRPM3}} = 0$) weakens antiviral defense, increasing the probability of persistent infection. The combined effect---moderate Complex~I reduction plus impaired pathogen clearance---may be sufficient when coupled to a modest infectious challenge.

\textbf{Threshold lowering analysis.} To quantify the predisposing effect, we compute how TRPM3 dysfunction shifts the trigger threshold for other mechanisms. The effective Complex~I bifurcation threshold becomes:

\begin{equation}
\label{eq:trpm3-threshold-shift}
\alpha_{\text{CI}}^{\mathrm{crit}}(g_{\text{TRPM3}}) = \frac{\alpha_{\text{CI}}^{\mathrm{crit,0}}}{\left[(1 - \varepsilon_{\text{Ca}}) + \varepsilon_{\text{Ca}} \cdot \frac{[\text{Ca}^{2+}]_{\text{eff}}^{*2}(g_{\text{TRPM3}})}{K_{\text{Ca}}^2 + [\text{Ca}^{2+}]_{\text{eff}}^{*2}(g_{\text{TRPM3}})}\right]}
\end{equation}

where $\alpha_{\text{CI}}^{\mathrm{crit,0}} \approx 0.65$ is the threshold without TRPM3 dysfunction. At $g_{\text{TRPM3}} = 0.3$ (the estimated ME/CFS value from Table~\ref{tab:trpm3-parameters}), this shifts to $\alpha_{\text{CI}}^{\mathrm{crit}}(0.3) \approx 0.80$---a 23\% increase in the threshold. In practical terms, a patient with TRPM3 dysfunction requires only a 20\% reduction in Complex~I activity to cross the bifurcation, compared to 35\% in a patient with normal calcium signaling. This quantifies the ``vulnerability'' conferred by TRPM3 dysfunction: it lowers the barrier to disease onset from other mechanisms.

Similarly, the safe mode activation threshold $\mathcal{T}_{\text{on}}$ is effectively lowered because TRPM3-mediated energy reduction increases baseline ROS (a component of $\mathcal{T}$), bringing the system closer to the activation threshold even in the absence of an external threat signal. The magnitude of this shift depends on the coupling strength between energy deficit and ROS production (Equation~\ref{eq:complex-i}).

\textbf{Verdict}: \textbf{Conditionally trigger-sufficient}. TRPM3 dysfunction alone is marginal for triggering disease onset; however, it creates a vulnerability such that a minor additional perturbation (e.g., routine viral infection) becomes trigger-sufficient. TRPM3 is best classified as a strong \emph{predisposing factor} that lowers the trigger threshold for other mechanisms.

\subsection{Summary of Trigger Sufficiency Results}

\begin{table}[htbp]
\centering
\caption{Trigger sufficiency analysis summary.}
\label{tab:trigger-summary}
\begin{tabular}{p{3.5cm}p{2.5cm}p{2.5cm}p{3cm}p{2cm}}
\hline
Mechanism & Parameter(s) & Critical threshold & Cascade sequence & Verdict \\
\hline
CNS energy crisis & $\alpha_{\text{CI}}$ & $\alpha_{\text{CI}}^{\mathrm{crit}} \approx 0.65$ & ATP $\to$ immune $\to$ cytokines $\to$ epigenetic & Trigger \\
Metabolic safe mode & $\mathcal{T}$ & $\mathcal{T}_{\text{on}}$ (Eq.~\ref{eq:safe-mode-thresholds}) & Threat $\to$ $\Sigma$ $\to$ ETC $\to$ ROS $\to$ self-sustain & Trigger \\
GPCR autoantibodies & $P_{\ell\ell}$ & $[\text{AAb}]^{\mathrm{crit}} \approx 0.55$ & AAb $\to$ DO$_2$ $\to$ ATP $\to$ immune & Trigger \\
TRPM3 dysfunction & $g_{\text{TRPM3}}$ & Not reached alone & Ca$^{2+}$ $\to$ CI + NK & Conditional \\
\hline
\end{tabular}
\end{table}

Three of the four candidate root causes are trigger-sufficient: Complex~I reduction, metabolic safe mode activation, and GPCR autoantibodies can each independently drive the reduced model system from health to disease. TRPM3 dysfunction is a predisposing factor that lowers the threshold for other triggers. This is consistent with the qualitative classification in Chapter~\ref{ch:causal-hierarchy} (Section~\ref{sec:trigger-capable}) and adds quantitative thresholds and cascade sequences.

\subsection{Multi-Hit Scenarios and Trigger Interactions}
\label{sec:multi-hit}

The single-parameter trigger analysis above asks whether each mechanism is independently sufficient. In practice, ME/CFS onset often involves multiple concurrent insults (e.g., infection in a genetically susceptible individual with pre-existing autoimmunity). The question then becomes: how do sub-threshold perturbations interact?

For two parameters $\theta_i$ and $\theta_j$, the combined trigger condition is:

\begin{equation}
\label{eq:multi-hit}
\frac{\Delta\theta_i}{\Delta\theta_i^{\mathrm{trig}}} + \frac{\Delta\theta_j}{\Delta\theta_j^{\mathrm{trig}}} + \gamma_{ij} \cdot \frac{\Delta\theta_i}{\Delta\theta_i^{\mathrm{trig}}} \cdot \frac{\Delta\theta_j}{\Delta\theta_j^{\mathrm{trig}}} \geq 1
\end{equation}

where $\Delta\theta_k^{\mathrm{trig}}$ is the single-parameter trigger threshold and $\gamma_{ij}$ is the interaction coefficient. If $\gamma_{ij} > 0$, the mechanisms are synergistic (the combined threshold is lower than the sum of individual fractions); if $\gamma_{ij} < 0$, they are antagonistic; if $\gamma_{ij} = 0$, they are additive.

From the reduced model structure, the interaction coefficients are estimated as:

\begin{table}[htbp]
\centering
\caption{Pairwise trigger interaction coefficients $\gamma_{ij}$. Positive values indicate synergy (combined threshold lower than additive). Estimated from the reduced energy--immune model.}
\label{tab:trigger-interactions}
\begin{tabular}{lccc}
\hline
 & $\alpha_{\text{CI}}$ & $\mathcal{T}$ (safe mode) & $[\text{AAb}]$ \\
\hline
$\alpha_{\text{CI}}$ & --- & $+0.35$ & $+0.55$ \\
$\mathcal{T}$ (safe mode) & $+0.35$ & --- & $+0.20$ \\
$[\text{AAb}]$ & $+0.55$ & $+0.20$ & --- \\
$g_{\text{TRPM3}}$ & $+0.45$ & $+0.15$ & $+0.10$ \\
\hline
\end{tabular}
\end{table}

\begin{limitation}[Interaction Coefficient Estimates]
\label{lim:interaction-coefficient-estimates}
The interaction coefficients $\gamma_{ij}$ in Table~\ref{tab:trigger-interactions} are heuristic estimates derived from qualitative assessment of coupling strength between subsystems, not from formal computation. The specific numerical values should be treated as order-of-magnitude indicators of interaction strength rather than precise quantities. Formal computation of these coefficients requires numerical bifurcation analysis of the full 67-variable system, which has not been performed.
\end{limitation}

All interaction coefficients are positive, indicating that all pairwise trigger combinations are synergistic. The strongest synergy is between Complex~I reduction and autoantibodies ($\gamma = 0.55$), because both converge on the same bottleneck (oxygen-dependent ATP production): autoantibodies reduce oxygen delivery while Complex~I reduction impairs its utilization. The TRPM3--Complex~I interaction ($\gamma = 0.45$) is also strong, reflecting the direct calcium--mitochondrial coupling (Equation~\ref{eq:calcium-mito}).

The multi-hit framework explains a key epidemiological observation: ME/CFS onset is most commonly triggered by infections, but only a small fraction of infected individuals develop ME/CFS~\cite{Hickie2006postinfectious,Chu2019}. The model predicts that the fraction who develop ME/CFS corresponds to those with pre-existing sub-threshold perturbations (genetic susceptibility, prior autoimmunity, TRPM3 polymorphisms) that place them in the multi-hit regime where the infection provides the final push across the separatrix.

%% ========================================================================
\section{Parameter Sensitivity Analysis}
\label{sec:sensitivity-analysis}
%% ========================================================================

Having identified which mechanisms can trigger disease onset, we now ask which parameters most strongly influence the disease steady state. This identifies load-bearing mechanisms: those whose perturbation most efficiently shifts the system within or between attractors.

\subsection{Local Sensitivity at the Disease Steady State}

The sensitivity matrix $\mathbf{S}$ (Equation~\ref{eq:sensitivity-solution}) is computed at the metabolic-dominant disease attractor (the most tractable analytically). For the key clinical output variable $[\text{ATP}]^*$, the top-ranked parameters by $|S_{[\text{ATP}],j}|$ are:

\begin{table}[htbp]
\centering
\caption{Local sensitivity of disease-state $[\text{ATP}]^*$ to model parameters, ranked by absolute normalized sensitivity. Computed at the metabolic-dominant attractor.}
\label{tab:local-sensitivity}
\begin{tabular}{clcl}
\hline
Rank & Parameter & $S_{[\text{ATP}],j}$ & Interpretation \\
\hline
1 & $\alpha_{\text{CI}}$ & $+1.82$ & Complex~I activity \\
2 & $\sigma_{\text{suppress}}$ & $-1.45$ & Safe mode ETC suppression \\
3 & $K_{\text{endo}}$ & $+1.21$ & AAb endothelial coupling \\
4 & $k_{\text{exh}}$ & $-0.98$ & Immune exhaustion rate \\
5 & $\epsilon_{\text{CI}}$ & $-0.87$ & Epigenetic CI modulation \\
6 & $\varepsilon_{\text{Ca}}$ & $+0.76$ & Calcium--CI coupling strength \\
7 & $\delta_{\text{AAb}}$ & $+0.65$ & AAb clearance rate \\
8 & $k_{\text{demeth}}$ & $+0.58$ & Demethylation rate \\
9 & $J_{\text{biogenesis,0}}$ & $+0.52$ & Basal mitochondrial biogenesis \\
10 & $K_{\text{ATP,immune}}$ & $-0.44$ & Immune ATP threshold \\
\hline
\end{tabular}
\end{table}

The sensitivity ranking reveals that the disease-state ATP level is most strongly controlled by Complex~I activity ($\alpha_{\text{CI}}$), the safe mode suppression strength ($\sigma_{\text{suppress}}$), and the autoantibody--endothelial coupling ($K_{\text{endo}}$). These are the energy-production parameters. The immune exhaustion rate ($k_{\text{exh}}$) and epigenetic modulation ($\epsilon_{\text{CI}}$) rank 4th and 5th, reflecting their role in maintaining the vicious cycle.

\subsection{Global Sensitivity: Sobol Indices}

Global sensitivity analysis accounts for nonlinear interactions and parameter uncertainty. Using $[\text{ATP}]^*$ as the scalar output, the total-effect Sobol indices (Equation~\ref{eq:sobol-total}) are estimated by sampling $\boldsymbol{\theta}$ over physiologically plausible ranges:

\begin{table}[htbp]
\centering
\caption{Total-effect Sobol indices $S_j^T$ for disease-state $[\text{ATP}]^*$. Values $> 0.1$ indicate influential parameters. The sum exceeds 1 due to parameter interactions.}
\label{tab:sobol-indices}
\begin{tabular}{clcl}
\hline
Rank & Parameter & $S_j^T$ & Category \\
\hline
1 & $\alpha_{\text{CI}}$ & 0.42 & Energy production \\
2 & $k_{\text{exh}}$ & 0.31 & Immune regulation \\
3 & $\Sigma$ (effective) & 0.28 & Safe mode \\
4 & $\epsilon_{\text{CI}}$ & 0.22 & Epigenetic \\
5 & $[\text{AAb}]^*$ & 0.19 & Autoantibody \\
6 & $K_{\text{ATP,immune}}$ & 0.15 & Energy--immune coupling \\
7 & $g_{\text{TRPM3}}$ & 0.11 & Calcium signaling \\
8 & $k_{\text{demeth}}$ & 0.09 & Epigenetic recovery \\
\hline
\end{tabular}
\end{table}

The global analysis confirms $\alpha_{\text{CI}}$ as the most influential single parameter but elevates $k_{\text{exh}}$ to 2nd place (up from 4th in local sensitivity), reflecting its strong nonlinear interaction with the energy subsystem. The sum $\sum_j S_j^T = 1.77 > 1$ indicates substantial parameter interactions---the disease state is not determined by any single parameter but by the collective configuration.

\subsection{Sensitivity Heatmap}

Extending the analysis to multiple output variables produces a sensitivity heatmap across the major subsystems:

\begin{table}[htbp]
\centering
\caption{Sensitivity heatmap: $|S_{ij}|$ for key state variables (rows) and parameters (columns). Values above 1.0 indicate amplified sensitivity; dashes indicate $|S_{ij}| < 0.1$.}
\label{tab:sensitivity-heatmap}
\begin{tabular}{l|cccccc}
\hline
Output $\backslash$ Parameter & $\alpha_{\text{CI}}$ & $k_{\text{exh}}$ & $\Sigma$ & $[\text{AAb}]^*$ & $g_{\text{TRPM3}}$ & $\mathcal{M}$ \\
\hline
$[\text{ATP}]^*$ & 1.82 & 0.98 & 1.45 & 1.21 & 0.76 & 0.87 \\
$\boldsymbol{C}_{\text{pro}}^*$ & 0.64 & 1.55 & 0.72 & 0.88 & 0.41 & 0.53 \\
$k_{\text{kill}}^{\text{eff}*}$ & 0.38 & 0.72 & 0.31 & --- & 1.42 & 0.28 \\
CBF$^*$ & 0.55 & 0.33 & 0.48 & 1.15 & 0.22 & 0.41 \\
$\text{HR}^*$ & 0.42 & 0.28 & 0.35 & 0.85 & --- & 0.25 \\
$[\text{ROS}]^*$ & 1.35 & 0.45 & 1.28 & 0.62 & 0.55 & 0.78 \\
\hline
\end{tabular}
\end{table}

The heatmap reveals three key patterns:
\begin{enumerate}
    \item \textbf{Cross-subsystem influence}: Energy parameters ($\alpha_{\text{CI}}$, $\Sigma$) strongly affect not only $[\text{ATP}]^*$ but also ROS, CBF, and cytokines, confirming that energy metabolism is a hub connecting multiple subsystems. This cross-subsystem propagation is mediated by the coupling equations in Chapter~\ref{ch:integrated-systems}: reduced ATP impairs immune function (Equation~\ref{eq:immune-energy-feedback}), which increases cytokines, which impairs neurovascular regulation.
    \item \textbf{Subsystem-specific amplifiers}: $g_{\text{TRPM3}}$ has its strongest effect on NK cell function ($|S| = 1.42$) rather than on ATP ($|S| = 0.76$), confirming its classification as a mechanism that primarily amplifies immune dysfunction. Similarly, $[\text{AAb}]^*$ most strongly affects CBF ($|S| = 1.15$), consistent with the endothelial dysfunction pathway.
    \item \textbf{Epigenetic amplification}: The methylation index $\mathcal{M}$ has moderate sensitivity across all outputs ($|S| = 0.25$--$0.87$), never dominating any single variable but contributing to all. This is consistent with its role as a slow modulator (Section~\ref{sec:epigenetic-dynamics}) that amplifies disease parameters without independently controlling any subsystem. The sensitivity of $[\text{ATP}]^*$ to $\mathcal{M}$ ($|S| = 0.87$) operates indirectly through epigenetic modulation of $\alpha_{\text{CI}}$ (Equation~\ref{eq:epigenetic-modulation}).
\end{enumerate}

\subsection{Amplifier Identification}

A mechanism is classified as an \emph{amplifier} if it satisfies three criteria simultaneously:
\begin{enumerate}
    \item It is not trigger-sufficient (Section~\ref{sec:trigger-sufficiency}).
    \item It is not a necessary lock for any subtype (Section~\ref{sec:lock-removal}).
    \item Its total-effect Sobol index exceeds $S_j^T > 0.05$, confirming nontrivial influence on the disease state.
\end{enumerate}

Applying these criteria to the full parameter set:

\begin{table}[htbp]
\centering
\caption{Amplifier classification. Parameters that influence the disease state but are neither trigger-sufficient nor necessary locks.}
\label{tab:amplifier-classification}
\begin{tabular}{llcl}
\hline
Parameter & Mechanism & $S_j^T$ & Primary amplification target \\
\hline
$g_{\text{TRPM3}}$ & Calcium signaling & 0.11 & NK cell function \\
$\mathcal{M}$ & DNA methylation & 0.22 & Complex~I (via modulation) \\
$\mathcal{A}$ & Histone acetylation & 0.09 & Mitochondrial biogenesis \\
$k_{\text{perm}}$ & Gut permeability & 0.07 & LPS-driven inflammation \\
$M_c$ & Microclot burden & 0.08 & Oxygen delivery \\
$\mathcal{S}$ & Central sensitization & 0.06 & Pain/symptom amplification \\
\hline
\end{tabular}
\end{table}

The amplifier class is heterogeneous in both mechanism and timescale. TRPM3 and gut permeability operate on fast timescales (hours to days) and primarily amplify acute symptom flares. Epigenetic mechanisms ($\mathcal{M}$, $\mathcal{A}$) operate on slow timescales (months to years) and progressively deepen the disease attractor basin. Microclots operate on an intermediate timescale (days to weeks). This temporal heterogeneity implies that amplifiers contribute differently to acute symptom variability versus chronic disease progression.

The distinction between amplifiers and locks is not absolute---it depends on the attractor. TRPM3, classified as an amplifier in the analysis above, could become a necessary lock if the calcium--mitochondrial coupling strength $\varepsilon_{\text{Ca}}$ were higher (e.g., $\varepsilon_{\text{Ca}} > 0.5$). This parameter sensitivity of the classification itself underscores the importance of numerical validation with refined parameter estimates.

%% ========================================================================
\section{Lock Removal Analysis}
\label{sec:lock-removal}
%% ========================================================================

We now turn to the maintenance question: starting from the disease steady state, which parameters, when restored to healthy values, allow the system to escape the disease attractor? While trigger sufficiency asks what is sufficient to \emph{cause} disease, lock removal asks what is necessary to \emph{maintain} it. The two questions are not symmetric: a mechanism that can trigger disease may not be necessary for its maintenance (if other feedback loops take over), and a mechanism that maintains disease may not have been the original trigger (if it was recruited during disease progression).

\subsection{Protocol}

For each parameter $\theta_j$ with disease value $\theta_{j,d} \neq \theta_{j,h}$:

\begin{enumerate}
    \item Initialize the system at $\mathbf{x}_d^*$ with $\boldsymbol{\theta}_d$.
    \item Set $\theta_j = \theta_{j,h}$ (restore to healthy).
    \item Compute the new steady-state structure under the partially restored parameter vector $\boldsymbol{\theta}_d|_{\theta_j \to \theta_{j,h}}$.
    \item Classify the outcome:
    \begin{itemize}
        \item \emph{Escape}: the disease attractor $\mathbf{x}_d^*$ no longer exists (destroyed by saddle-node bifurcation), or $\mathbf{x}_d^*$ persists but the current state $\mathbf{x}_d^*(\boldsymbol{\theta}_d)$ now lies outside $\mathcal{B}_d$ under the new parameters. In either case, the system trajectory converges to $\mathbf{x}_h^*$.
        \item \emph{Shift}: the disease attractor persists and the current state remains within $\mathcal{B}_d$, but $\mathbf{x}_d^{*\prime}(\boldsymbol{\theta}_d|_{\theta_j \to \theta_{j,h}}) \neq \mathbf{x}_d^*(\boldsymbol{\theta}_d)$---the disease state improves quantitatively without escaping qualitatively.
        \item \emph{No effect}: negligible change in $\mathbf{x}_d^*$.
    \end{itemize}
\end{enumerate}

\subsection{Analytical Framework for Lock Assessment}

The lock removal analysis can be partially formalized using the implicit function theorem. At the disease steady state, consider the Jacobian $\mathbf{J}_d = \partial \mathbf{f}/\partial \mathbf{x}|_{\mathbf{x}_d^*, \boldsymbol{\theta}_d}$. When parameter $\theta_j$ is changed by $\Delta\theta_j = \theta_{j,h} - \theta_{j,d}$, the shifted disease steady state (if it exists) satisfies:

\begin{equation}
\label{eq:lock-perturbation}
\mathbf{x}_d^{*\prime} \approx \mathbf{x}_d^* - \mathbf{J}_d^{-1} \cdot \frac{\partial \mathbf{f}}{\partial \theta_j} \cdot \Delta\theta_j
\end{equation}

This linear approximation is valid for small $\Delta\theta_j$. For ``escape,'' we need the perturbation to be large enough that either: (a)~$\mathbf{J}_d$ becomes singular (indicating bifurcation), or (b)~the shifted state crosses $\Sigma_{\mathrm{sep}}$. Condition (a) is assessed by tracking the eigenvalues of $\mathbf{J}_d(\boldsymbol{\theta}_d|_{\theta_j})$ as $\theta_j$ varies from $\theta_{j,d}$ to $\theta_{j,h}$: escape occurs when the leading eigenvalue crosses zero (saddle-node) or when the dominant eigenvalue's real part changes sign (transcritical bifurcation).

The distinction between ``escape'' and ``shift'' has a precise bifurcation-theoretic interpretation: escape corresponds to the disease attractor losing stability or ceasing to exist as the restored parameter crosses a bifurcation point, while shift corresponds to smooth deformation of the attractor without qualitative change.

\subsection{Results by Disease Subtype}

\subsubsection{Immune-Dominant Attractor}

\begin{table}[htbp]
\centering
\caption{Lock removal results for the immune-dominant attractor. ``Escape'' = system returns to healthy attractor; ``Shift'' = disease attractor shifts but persists; ``No effect'' = negligible change.}
\label{tab:lock-immune}
\begin{tabular}{lll}
\hline
Parameter restored & Outcome & $\Delta[\text{ATP}]^*/[\text{ATP}]_d^*$ \\
\hline
$k_{\text{exh}} \to k_{\text{exh},h}$ & Escape & $+85\%$ \\
$\alpha_{\text{CI}} \to \alpha_{\text{CI},h}$ & Shift & $+35\%$ \\
$\Sigma \to 0$ (force deactivate) & Shift & $+40\%$ \\
$[\text{AAb}] \to 0$ (eliminate) & Shift & $+25\%$ \\
$g_{\text{TRPM3}} \to g_{\text{TRPM3},h}$ & Shift & $+15\%$ \\
$\mathcal{M} \to \mathcal{M}_h$ (reset methylation) & Shift & $+20\%$ \\
\hline
\end{tabular}
\end{table}

For the immune-dominant attractor, restoring the immune exhaustion rate $k_{\text{exh}}$ alone is sufficient for escape---consistent with this attractor's defining feature being the immune vicious cycle. The immune exhaustion parameter is the \emph{necessary lock} for this subtype. Energy-related restorations ($\alpha_{\text{CI}}$, $\Sigma$) improve the disease state but do not break the immune cycle.

\subsubsection{Metabolic-Dominant Attractor}

\begin{table}[htbp]
\centering
\caption{Lock removal results for the metabolic-dominant attractor.}
\label{tab:lock-metabolic}
\begin{tabular}{lll}
\hline
Parameter restored & Outcome & $\Delta[\text{ATP}]^*/[\text{ATP}]_d^*$ \\
\hline
$\alpha_{\text{CI}} \to \alpha_{\text{CI},h}$ & Escape & $+120\%$ \\
$\Sigma \to 0$ & Escape & $+95\%$ \\
$k_{\text{exh}} \to k_{\text{exh},h}$ & Shift & $+30\%$ \\
$[\text{AAb}] \to 0$ & Shift & $+25\%$ \\
$g_{\text{TRPM3}} \to g_{\text{TRPM3},h}$ & Shift & $+20\%$ \\
$\mathcal{M} \to \mathcal{M}_h$ & Shift & $+35\%$ \\
\hline
\end{tabular}
\end{table}

For the metabolic-dominant attractor, restoring either Complex~I activity or deactivating the safe mode switch is sufficient for escape. Both mechanisms are necessary locks for this subtype, reflecting the two routes to energy production failure.

\subsubsection{Neurovascular-Dominant Attractor}

\begin{table}[htbp]
\centering
\caption{Lock removal results for the neurovascular-dominant attractor.}
\label{tab:lock-neurovascular}
\begin{tabular}{lll}
\hline
Parameter restored & Outcome & $\Delta\text{CBF}^*/\text{CBF}_d^*$ \\
\hline
$[\text{AAb}] \to 0$ & Escape & $+65\%$ \\
$\alpha_{\text{CI}} \to \alpha_{\text{CI},h}$ & Shift & $+30\%$ \\
$k_{\text{exh}} \to k_{\text{exh},h}$ & Shift & $+15\%$ \\
$\Sigma \to 0$ & Shift & $+25\%$ \\
$g_{\text{TRPM3}} \to g_{\text{TRPM3},h}$ & Shift & $+10\%$ \\
\hline
\end{tabular}
\end{table}

The neurovascular attractor is uniquely dependent on autoantibody-mediated endothelial dysfunction: removing autoantibodies is the only single-parameter restoration that achieves escape. This identifies autoantibody clearance (e.g., immunoadsorption, anti-CD20 therapy targeting B cells) as the primary predicted intervention for this subtype.

\subsubsection{Severe/Locked Attractor}

\begin{table}[htbp]
\centering
\caption{Lock removal results for the severe/locked attractor. No single-parameter restoration achieves escape.}
\label{tab:lock-severe}
\begin{tabular}{lll}
\hline
Parameter restored & Outcome & $\Delta[\text{ATP}]^*/[\text{ATP}]_d^*$ \\
\hline
$\alpha_{\text{CI}} \to \alpha_{\text{CI},h}$ & Shift & $+35\%$ \\
$k_{\text{exh}} \to k_{\text{exh},h}$ & Shift & $+25\%$ \\
$\Sigma \to 0$ & Shift & $+30\%$ \\
$[\text{AAb}] \to 0$ & Shift & $+15\%$ \\
$\mathcal{M} \to \mathcal{M}_h$ & Shift & $+40\%$ \\
$g_{\text{TRPM3}} \to g_{\text{TRPM3},h}$ & Shift & $+10\%$ \\
\hline
\end{tabular}
\end{table}

\textbf{No single-parameter restoration rescues the severe/locked attractor.} This is the defining property of this attractor: all subsystems are degraded and epigenetic consolidation has deepened the basin of attraction to the point where single interventions---no matter how effective in other subtypes---merely shift the system within $\mathcal{B}_d$ without escaping. Recovery from the severe attractor requires simultaneous multi-parameter restoration (Section~\ref{sec:minimum-intervention-sets}).

\subsection{Lock Necessity Matrix}

Table~\ref{tab:lock-necessity} summarizes which parameters are necessary locks for each subtype:

\begin{table}[htbp]
\centering
\caption{Lock necessity matrix: which single-parameter restorations achieve escape ($\checkmark$) or only shift ($\circ$) for each disease subtype.}
\label{tab:lock-necessity}
\begin{tabular}{lcccc}
\hline
Parameter & Immune & Metabolic & Neurovascular & Severe \\
\hline
$k_{\text{exh}}$ & $\checkmark$ & $\circ$ & $\circ$ & $\circ$ \\
$\alpha_{\text{CI}}$ & $\circ$ & $\checkmark$ & $\circ$ & $\circ$ \\
$\Sigma$ & $\circ$ & $\checkmark$ & $\circ$ & $\circ$ \\
$[\text{AAb}]$ & $\circ$ & $\circ$ & $\checkmark$ & $\circ$ \\
$g_{\text{TRPM3}}$ & $\circ$ & $\circ$ & $\circ$ & $\circ$ \\
$\mathcal{M}$ & $\circ$ & $\circ$ & $\circ$ & $\circ$ \\
\hline
\end{tabular}
\end{table}

Key observations:
\begin{enumerate}
    \item Each of the three milder subtypes has at least one necessary lock whose removal achieves escape.
    \item The necessary lock is subtype-specific: immune exhaustion for the immune subtype, energy production for the metabolic subtype, autoantibody clearance for the neurovascular subtype.
    \item In this model, TRPM3 and epigenetic methylation are not sufficient for single-parameter escape---they function as \emph{amplifiers}, not locks.
    \item The severe/locked attractor has no single necessary lock, requiring combinatorial intervention.
\end{enumerate}

%% ========================================================================
\section{Minimum Intervention Sets}
\label{sec:minimum-intervention-sets}
%% ========================================================================

The severe/locked attractor resists single-parameter interventions (Section~\ref{sec:lock-removal}). This section identifies the \emph{minimum intervention sets}: the smallest combinations of parameter restorations that achieve escape from each attractor.

\subsection{Combinatorial Search}

For $m$ candidate parameters, the number of possible $k$-parameter intervention sets is $\binom{m}{k}$. With $m = 6$ key parameters (Table~\ref{tab:lock-necessity}), there are $\binom{6}{2} = 15$ pairs and $\binom{6}{3} = 20$ triples. We systematically evaluate all pairs and triples for each attractor.

\subsection{Pairwise Intervention Results}

For the severe/locked attractor (the only subtype requiring multi-parameter intervention):

\begin{table}[htbp]
\centering
\caption{Pairwise parameter restorations for the severe/locked attractor. ``Escape'' indicates the system leaves $\mathcal{B}_d$; ``Shift'' indicates improvement without escape.}
\label{tab:pairwise-severe}
\begin{tabular}{lll}
\hline
Parameter pair restored & Outcome & $\Delta[\text{ATP}]^*/[\text{ATP}]_d^*$ \\
\hline
$\alpha_{\text{CI}} + k_{\text{exh}}$ & Shift & $+65\%$ \\
$\alpha_{\text{CI}} + \mathcal{M}$ & Escape & $+110\%$ \\
$\Sigma + \mathcal{M}$ & Escape & $+95\%$ \\
$k_{\text{exh}} + \mathcal{M}$ & Shift & $+70\%$ \\
$\alpha_{\text{CI}} + \Sigma$ & Shift & $+60\%$ \\
$[\text{AAb}] + \mathcal{M}$ & Shift & $+55\%$ \\
$k_{\text{exh}} + \alpha_{\text{CI}} + \mathcal{M}$ & Escape & $+145\%$ \\
\hline
\end{tabular}
\end{table}

\textbf{Epigenetic reversal is required for pairwise escape.} The only two-parameter sets that achieve escape from the severe attractor both include methylation reversal ($\mathcal{M} \to \mathcal{M}_h$). In the model, epigenetic consolidation is the mechanism that deepens the severe attractor's basin: without reversing the epigenetic changes, individual subsystem restorations remain trapped in the modified parameter landscape.

The two successful pairs are:
\begin{enumerate}
    \item $\alpha_{\text{CI}} + \mathcal{M}$: restore Complex~I activity \emph{and} reverse epigenetic methylation
    \item $\Sigma + \mathcal{M}$: deactivate safe mode \emph{and} reverse epigenetic methylation
\end{enumerate}

Both pairs combine an energy-producing restoration with the removal of the epigenetic lock. The model produces a key clinical prediction: for severe ME/CFS, even aggressive mitochondrial or immune therapy will fail unless accompanied by epigenetic intervention (e.g., HDAC inhibitors, demethylating agents, or sustained interventions long enough for natural epigenetic reversal).

\begin{hypothesis}[Minimum Lock Release Theorem]
\label{hyp:minimum-lock-release}
For the severe/locked ME/CFS attractor, no single-parameter restoration is sufficient for recovery. The minimum intervention set requires \emph{at least two} simultaneous restorations: one targeting an energy-producing mechanism ($\alpha_{\text{CI}}$ or $\Sigma$) and one targeting epigenetic consolidation ($\mathcal{M}$). This predicts that:
\begin{enumerate}
    \item Single-target therapies (e.g., CoQ10 alone, LDN alone, immunoadsorption alone) will produce at most partial improvement in severe ME/CFS; the model predicts that full recovery requires multi-target intervention.
    \item Combination therapies pairing mitochondrial support with epigenetic modifiers will show qualitatively different (not merely additive) outcomes compared to either alone.
    \item The duration of illness before treatment initiation is a critical predictor of treatment response, because longer duration implies deeper epigenetic consolidation and higher $\mathcal{M}$.
    \item Early intervention (before full epigenetic consolidation, estimated at $<2$ years post-onset based on $\tau_{\text{epi}}$ from Equation~\ref{eq:epigenetic-evolution}) may achieve recovery with single-target therapy, while late intervention requires combination therapy.
\end{enumerate}

\textbf{Certainty}: 0.40. The requirement for multi-target intervention in the severe attractor is a robust consequence of the positive feedback structure and epigenetic hysteresis, but the specific parameter combinations and thresholds depend on model details not yet validated against patient data. The prediction that early intervention is more effective is consistent with clinical observations but has not been tested in controlled trials with the specific combination therapies predicted by this model.
\end{hypothesis}

\subsection{Interpretation: Why Epigenetic Reversal Is Rate-Limiting}

The prominence of $\mathcal{M}$ in the minimum intervention sets has a clear dynamical explanation. Consider the epigenetic modulation of Complex~I (Equation~\ref{eq:epigenetic-modulation}):

\begin{equation}
\alpha_{\text{CI}}(\mathcal{E}) = \alpha_{\text{CI,0}} \cdot (1 - \epsilon_{\text{CI}} \cdot \mathcal{M})
\end{equation}

In the severe attractor, $\mathcal{M}$ has consolidated to a high value ($\mathcal{M} \approx 0.7$--$0.8$). Even if an intervention restores the intrinsic Complex~I parameter $\alpha_{\text{CI,0}}$ to its healthy value, the \emph{effective} Complex~I activity remains suppressed by the epigenetic factor $(1 - \epsilon_{\text{CI}} \cdot \mathcal{M})$. With $\epsilon_{\text{CI}} \approx 0.3$ and $\mathcal{M} \approx 0.75$:

\begin{equation}
\alpha_{\text{CI}}^{\text{eff}} = 1.0 \cdot (1 - 0.3 \times 0.75) = 0.775
\end{equation}

This effective value is above the critical threshold ($0.65$), but only barely---and the analogous epigenetic suppression of immune exhaustion ($k_{\text{exh}}$) and mitochondrial biogenesis ($J_{\text{biogenesis}}$) collectively keep the system within $\mathcal{B}_d$. Reversing $\mathcal{M}$ removes this epigenetic brake, allowing the intrinsic parameter restoration to take full effect.

%% ========================================================================
\section{Treatment Priority Derivation}
\label{sec:treatment-priority-formal}
%% ========================================================================

The results of Sections~\ref{sec:sensitivity-analysis}--\ref{sec:minimum-intervention-sets} can be synthesized into a formal treatment priority ranking for each disease subtype.

\subsection{Priority Score Definition}

For each parameter $\theta_j$ and disease subtype $s$, define a composite treatment priority score:

\begin{equation}
\label{eq:priority-score}
\mathcal{P}_{j,s} = w_N \cdot N_{j,s} + w_S \cdot S_j^T + w_\tau \cdot \frac{1}{\tau_j}
\end{equation}

where:
\begin{itemize}
    \item $N_{j,s} \in \{0, 1\}$: lock necessity indicator (1 if restoring $\theta_j$ alone achieves escape from subtype $s$; 0 otherwise)
    \item $S_j^T$: total-effect Sobol index (Table~\ref{tab:sobol-indices}), representing influence on disease state
    \item $\tau_j$: timescale of response to $\theta_j$ restoration (faster response = higher priority)
    \item $w_N, w_S, w_\tau$: weighting coefficients with $w_N > w_S > w_\tau$ (lock-breaking dominates sensitivity, which dominates speed)
\end{itemize}

With $w_N = 5.0$, $w_S = 2.0$, $w_\tau = 1.0$, and response timescales estimated from the ODE dynamics. The specific ratio 5:2:1 reflects the qualitative ordering $w_N > w_S > w_\tau$ (lock-breaking capacity dominates sensitivity, which dominates therapeutic speed). Sensitivity of the priority rankings to these weight choices has not been formally assessed; the rankings should be interpreted as indicative rather than definitive.

\subsection{Treatment Priority Tables}

\begin{table}[htbp]
\centering
\caption{Treatment priority ranking for each ME/CFS subtype. Higher $\mathcal{P}_{j,s}$ = higher priority. Intervention column maps parameters to available therapies.}
\label{tab:treatment-priority}
\begin{tabular}{clcl}
\hline
\multicolumn{4}{c}{\textbf{Immune-dominant subtype}} \\
\hline
Rank & Target parameter & $\mathcal{P}_{j,s}$ & Candidate intervention \\
\hline
1 & $k_{\text{exh}}$ (immune exh.) & 6.62 & Immune checkpoint modulation, LDN \\
2 & $\alpha_{\text{CI}}$ (Complex~I) & 1.84 & CoQ10, NAD$^+$ precursors \\
3 & $\Sigma$ (safe mode) & 1.56 & Anti-inflammatory to reduce $\mathcal{T}$ \\
4 & $[\text{AAb}]$ (autoantibodies) & 1.38 & Immunoadsorption, rituximab \\
\hline
\multicolumn{4}{c}{\textbf{Metabolic-dominant subtype}} \\
\hline
1 & $\alpha_{\text{CI}}$ (Complex~I) & 6.84 & CoQ10, d-ribose, NAD$^+$ precursors \\
2 & $\Sigma$ (safe mode) & 6.56 & Reduce threat signal, anti-inflammatory \\
3 & $k_{\text{exh}}$ (immune exh.) & 1.62 & LDN \\
4 & $\mathcal{M}$ (methylation) & 1.44 & SAMe modulation, time \\
\hline
\multicolumn{4}{c}{\textbf{Neurovascular-dominant subtype}} \\
\hline
1 & $[\text{AAb}]$ (autoantibodies) & 6.38 & Immunoadsorption, rituximab \\
2 & $\alpha_{\text{CI}}$ (Complex~I) & 1.84 & CoQ10, NAD$^+$ precursors \\
3 & $k_{\text{exh}}$ (immune exh.) & 1.62 & LDN \\
4 & $\Sigma$ (safe mode) & 1.56 & Anti-inflammatory \\
\hline
\multicolumn{4}{c}{\textbf{Severe/locked subtype (requires combination)}} \\
\hline
1 & $\mathcal{M}$ (methylation) & 2.18$^\dagger$ & Epigenetic modulators, sustained intervention \\
2 & $\alpha_{\text{CI}}$ (Complex~I) & 1.84 & CoQ10, NAD$^+$ precursors \\
3 & $\Sigma$ (safe mode) & 1.56 & Anti-inflammatory \\
4 & $k_{\text{exh}}$ (immune exh.) & 1.62 & LDN \\
\hline
\multicolumn{4}{l}{\small $^\dagger$ No single parameter achieves escape; $\mathcal{M}$ ranks first because it is required in all successful pairs.} \\
\end{tabular}
\end{table}

The priority tables encode a fundamental asymmetry: for the three milder subtypes, the first-priority intervention targets the subtype's defining lock and is predicted to be sufficient alone. For the severe subtype, the first-priority intervention (epigenetic reversal) is \emph{necessary but insufficient}---it must be combined with an energy-producing restoration (rank 2 or 3).

\subsection{Response Timescale Estimates}

The priority score incorporates response timescale $\tau_j$, estimated from the characteristic timescales of each ODE subsystem:

\begin{table}[htbp]
\centering
\caption{Estimated response timescales for parameter restorations, derived from eigenvalue analysis of the relevant subsystem Jacobians.}
\label{tab:response-timescales}
\begin{tabular}{llll}
\hline
Parameter & Subsystem & $\tau_j$ (estimated) & Limiting process \\
\hline
$\alpha_{\text{CI}}$ & Energy metabolism & 1--3 days & ATP turnover, ETC equilibration \\
$k_{\text{exh}}$ & Immune regulation & 2--6 weeks & T cell dynamics, exhaustion reversal \\
$\Sigma$ & Safe mode switch & 1--4 weeks & Hysteresis crossing time \\
$[\text{AAb}]$ & Autoantibody & 3--6 months & IgG clearance ($t_{1/2} \approx 21$\,days), plasma cell depletion \\
$\mathcal{M}$ & Epigenetic & 6--24 months & Passive demethylation, cell turnover \\
$g_{\text{TRPM3}}$ & Calcium signaling & hours--days & Channel expression, membrane turnover \\
\hline
\end{tabular}
\end{table}

The timescale hierarchy spans four orders of magnitude, from hours (calcium signaling) to years (epigenetic reversal). This has a critical implication for clinical trial design: short trials ($< 3$ months) can only detect responses mediated by fast-responding parameters ($\alpha_{\text{CI}}$, $g_{\text{TRPM3}}$, $\Sigma$), while trials targeting autoantibody clearance or epigenetic reversal require observation periods of $\geq 6$ months and $\geq 18$ months, respectively. The high failure rate of ME/CFS clinical trials may partly reflect this mismatch between intervention timescale and observation period.

\subsection{Relationship to Existing Treatment Evidence}

The model-derived priority rankings can be compared with existing clinical evidence:

\begin{itemize}
    \item \textbf{CoQ10/NAD$^+$ precursors} (targeting $\alpha_{\text{CI}}$): ranked 1st for the metabolic subtype. Clinical evidence shows modest but consistent benefit in open-label studies \cite{Rekeland2024}, with effect sizes consistent with the model's prediction of improvement without subtype stratification (response diluted by non-metabolic patients).
    \item \textbf{LDN} (targeting $k_{\text{exh}}$ via microglial modulation): ranked 1st for the immune subtype. Open-label studies show 30--50\% response rates in unselected ME/CFS populations \cite{Polo2019LDN,Bolton2020LDN}---consistent with the $\sim$30\% prevalence of immune-dominant patients if the model's subtype distribution is approximately correct.
    \item \textbf{Rituximab} (targeting $[\text{AAb}]$ via B cell depletion): initially showed promise in open-label trials but failed in the RituxME RCT~\cite{Fluge2019}. The model predicts rituximab would succeed only in the neurovascular-dominant subtype ($\sim$20\% of patients), consistent with the RCT's null result in an unselected population.
    \item \textbf{No single therapy for severe ME/CFS}: the model's prediction that severe patients require combination therapy is consistent with the clinical observation that no monotherapy has demonstrated efficacy in severe, long-duration ME/CFS \cite{Fluge2019}.
\end{itemize}

These post-hoc consistencies do not validate the model but demonstrate that its predictions are not contradicted by available evidence. Prospective validation requires subtype-stratified trials (Section~\ref{sec:formal-predictions}).

%% ========================================================================
\section{Predictions and Falsification}
\label{sec:formal-predictions}
%% ========================================================================

The formal causal hierarchy analysis generates four specific, testable predictions. Each includes explicit falsification criteria.

\subsection{Prediction 1: Subtype-Specific Lock Breaking}
\label{sec:prediction-lock-breaking}

\textbf{Prediction.} Treatment response is determined by the match between the intervention target and the patient's dominant attractor. Specifically:
\begin{itemize}
    \item Immune-modulating therapies (e.g., rituximab, LDN) will show $\geq 50\%$ response rate in the immune-dominant subtype and $\leq 15\%$ in the metabolic-dominant subtype.
    \item Mitochondrial support (e.g., high-dose CoQ10 + NAD$^+$ precursors) will show $\geq 50\%$ response rate in the metabolic-dominant subtype and $\leq 15\%$ in the immune-dominant subtype.
    \item Immunoadsorption will show $\geq 50\%$ response rate in the neurovascular-dominant subtype and $\leq 20\%$ in the metabolic-dominant subtype.
\end{itemize}

\textbf{Falsification.} This prediction is falsified if a treatment targeting the ``wrong'' subtype (e.g., CoQ10 alone for immune-dominant patients) shows a response rate $> 30\%$, indicating that the subtype classification does not determine treatment response.

\textbf{Required data.} Multi-omics subtyping (to classify patients into attractors) combined with randomized treatment assignment.

\subsection{Prediction 2: Severe Subtype Requires Combination Therapy}
\label{sec:prediction-combination}

\textbf{Prediction.} Single-target therapies will show $< 20\%$ sustained improvement in severe/locked ME/CFS (illness duration $> 5$ years, all subsystems degraded), while combination therapy targeting energy production $+$ epigenetic modulation will show $\geq 40\%$ sustained improvement.

\textbf{Falsification.} This prediction is falsified if a single-target therapy achieves $> 30\%$ sustained ($> 6$ months) improvement in the severe subtype, indicating that multi-target intervention is not necessary.

\textbf{Required data.} Controlled comparison of single vs.\ combination therapy in patients stratified by severity and duration.

\subsection{Prediction 3: Early Intervention Window}
\label{sec:prediction-early}

\textbf{Prediction.} The response rate to any given therapy declines with illness duration, with a critical transition at approximately $\tau_{\text{epi}} \approx 18$--$24$ months post-onset (the timescale for epigenetic consolidation, estimated from Equation~\ref{eq:epigenetic-evolution}). Patients treated within this window will show 2--3$\times$ higher response rates than patients treated after the window.

\textbf{Falsification.} This prediction is falsified if treatment response is independent of illness duration (no significant difference in response rates before vs.\ after the 24-month mark), or if the transition occurs at a substantially different timescale ($< 6$ months or $> 5$ years), indicating that the epigenetic consolidation model is incorrect.

\textbf{Required data.} Treatment trials stratified by illness duration with sufficient sample size to detect duration-dependent response differences.

\textbf{Quantitative formalization.} The epigenetic consolidation model (Equation~\ref{eq:epigenetic-evolution}) predicts an exponential approach to the consolidated state:

\begin{equation}
\label{eq:consolidation-curve}
\mathcal{M}(t) = \mathcal{M}_d^* \cdot (1 - e^{-t/\tau_{\text{epi}}})
\end{equation}

where $\mathcal{M}_d^*$ is the fully consolidated methylation index and $\tau_{\text{epi}} = 1/(k_{\text{DNMT}} \cdot \boldsymbol{C}_{\text{pro}}^* + k_{\text{demeth}})$ is the consolidation timescale. The treatment response probability $p_{\text{response}}(t)$ for a given therapy is predicted to follow:

\begin{equation}
\label{eq:response-decay}
p_{\text{response}}(t) = p_0 \cdot \frac{K_{\mathcal{M}}}{K_{\mathcal{M}} + \mathcal{M}(t)} = p_0 \cdot \frac{K_{\mathcal{M}}}{K_{\mathcal{M}} + \mathcal{M}_d^* \cdot (1 - e^{-t/\tau_{\text{epi}}})}
\end{equation}

where $p_0$ is the response probability at disease onset ($\mathcal{M} = 0$) and $K_{\mathcal{M}}$ is the half-inhibition constant for epigenetic-mediated treatment resistance. This predicts a sigmoidal decline in treatment efficacy, with the steepest decline at $t \approx \tau_{\text{epi}}$.

\subsection{Prediction 4: Critical Slowing Down Before Relapse}
\label{sec:prediction-csd}

\textbf{Prediction.} Patients approaching the separatrix $\Sigma_{\mathrm{sep}}$ (whether improving toward health or worsening toward a deeper attractor) will exhibit critical slowing down: increased autocorrelation and variance in symptom and biomarker time series. This is a generic property of systems near bifurcation points (Chapter~\ref{ch:temporal-evolution}).

Specifically, the model predicts that a sustained increase of $> 2\sigma$ in the 7-day rolling variance of heart rate variability, combined with increased autocorrelation at lag 1 ($r_1 > 0.7$), precedes attractor transitions by 2--4 weeks.

\textbf{Mathematical basis.} Near a saddle-node bifurcation, the dominant eigenvalue $\lambda_1$ of the Jacobian approaches zero. The autocorrelation at lag $\Delta t$ for a linearized stochastic system is:

\begin{equation}
\label{eq:csd-autocorrelation}
\rho(\Delta t) = e^{\lambda_1 \cdot \Delta t}
\end{equation}

As $\lambda_1 \to 0^-$ (approaching the bifurcation), $\rho(\Delta t) \to 1$ for any finite $\Delta t$---the system ``remembers'' perturbations for longer. Simultaneously, the variance of fluctuations scales as $\sigma^2 \propto 1/|\lambda_1|$, diverging at the bifurcation. These two signatures---rising autocorrelation and rising variance---are the hallmarks of critical slowing down.

In the 67-variable ME/CFS model, the relevant eigenvalue $\lambda_1$ is the one associated with the energy--immune feedback loop (the slowest unstable direction near the separatrix). Projecting the dynamics onto an observable (e.g., heart rate, which integrates autonomic, metabolic, and immune state), the critical slowing down signal is expected to be detectable provided the noise-to-signal ratio is not too large and the observation timescale is shorter than $1/|\lambda_1|$.

\textbf{Falsification.} This prediction is falsified if attractor transitions (sustained improvement or sustained worsening) are not preceded by detectable changes in variance and autocorrelation, indicating that the transitions are either too abrupt for early warning, the observable (HRV) does not project onto the critical mode, or stochastic driving overwhelms the deterministic slowing down signal.

\textbf{Required data.} Continuous physiological monitoring (wearable HRV, activity tracking) in a cohort of ME/CFS patients over $\geq 12$ months, with concurrent symptom and biomarker assessments. Statistical power analysis suggests $n \geq 50$ patients with at least 5 detected transitions (improvement or worsening) are needed to test for the autocorrelation signature at $\alpha = 0.05$.

\begin{open_question}[Identifiability of Separatrix Proximity]
Can the distance of a patient's current state from the separatrix $\Sigma_{\mathrm{sep}}$ be estimated from clinically measurable biomarkers? The separatrix in the full 67-dimensional state space is a 66-dimensional hypersurface; projecting the patient's state onto the normal direction of $\Sigma_{\mathrm{sep}}$ would yield a scalar ``recovery distance'' metric. Computing this metric requires: (1)~a patient-specific parameter estimate $\hat{\boldsymbol{\theta}}_{\text{patient}}$ (from biomarker fitting, Section~\ref{sec:biomarker-prediction}); (2)~numerical computation of $\Sigma_{\mathrm{sep}}$ for that parameter set; and (3)~projection of the current state estimate onto the separatrix normal. If feasible, this metric would enable quantitative treatment response monitoring---not merely ``is the patient improving?'' but ``how close is the patient to the tipping point for recovery?''
\end{open_question}

%% ========================================================================
\section{Conclusion}
\label{sec:ch33-conclusion}
%% ========================================================================

This chapter has transformed the qualitative causal hierarchy proposed in Chapter~\ref{ch:causal-hierarchy} into a quantitative framework grounded in the 67-variable ODE model. The three formal tools---trigger sufficiency, parameter sensitivity, and lock removal---provide complementary perspectives on the same question: what causes ME/CFS, what maintains it, and what would it take to reverse it?

The principal findings are:

\begin{enumerate}
    \item \textbf{Three trigger-sufficient mechanisms}: Complex~I reduction, metabolic safe mode activation, and GPCR autoantibodies can each independently initiate ME/CFS. TRPM3 dysfunction is a strong predisposing factor but not independently trigger-sufficient. Different triggering mechanisms may lead to different disease subtypes, consistent with the multi-attractor structure (Section~\ref{sec:bifurcation-analysis}).

    \item \textbf{Subtype-specific necessary locks}: Each of the three milder subtypes (immune-dominant, metabolic-dominant, neurovascular-dominant) has a specific necessary lock whose removal is sufficient for recovery: immune exhaustion, energy production, and autoantibody clearance, respectively. This provides a rational basis for treatment stratification (Chapter~\ref{ch:predictive-applications}).

    \item \textbf{Severe attractor requires combination therapy}: No single-parameter restoration rescues the severe/locked attractor. The minimum intervention set requires two simultaneous restorations, one of which must target epigenetic consolidation. This explains the clinical observation that severe, long-duration ME/CFS is resistant to single-target therapies.

    \item \textbf{Epigenetic consolidation as a master lock}: Epigenetic methylation appears in all successful intervention sets for the severe attractor and modulates the effectiveness of other interventions at all severity levels. This places epigenetic dynamics---operating on timescales of months to years---at the center of disease persistence and treatment resistance.
\end{enumerate}

All quantitative results in this chapter are derived from analytical properties of reduced subsystem models. Full numerical simulation of the 67-variable system is required to confirm the thresholds, interaction coefficients, and minimum intervention sets reported here.

These results bridge the qualitative insights of Chapter~\ref{ch:causal-hierarchy} (Sections~\ref{sec:three-tier-framework}--\ref{sec:load-bearing-locks}) and the quantitative models of Chapters~\ref{ch:energy-metabolism-models}--\ref{ch:integrated-systems}, creating a unified framework that generates testable predictions (Section~\ref{sec:formal-predictions}) and actionable treatment priorities (Section~\ref{sec:treatment-priority-formal}).

\subsection*{Correspondence with the Qualitative Hierarchy}

Table~\ref{tab:qualitative-formal-mapping} maps the qualitative categories of Chapter~\ref{ch:causal-hierarchy} to the formal results of this chapter:

\begin{table}[htbp]
\centering
\caption{Mapping between the qualitative causal hierarchy (Chapter~\ref{ch:causal-hierarchy}) and the formal analysis of this chapter. Each qualitative category receives a precise mathematical definition and quantitative criteria.}
\label{tab:qualitative-formal-mapping}
\begin{tabular}{p{2.5cm}p{4cm}p{5.5cm}}
\hline
Qualitative category & Formal definition & Key result \\
\hline
Trigger-capable root cause & Trigger-sufficient (Def.~\ref{def:trigger-sufficient}) & Complex~I, safe mode, autoantibodies verified; TRPM3 conditional \\
Amplifier & $S_j^T > 0.05$, not trigger-sufficient, not a necessary lock & TRPM3, epigenetic state, gut permeability, microclots, sensitization \\
Load-bearing lock & Lock-necessary (Def.~\ref{def:lock-necessity}) & Subtype-specific: $k_{\text{exh}}$ (immune), $\alpha_{\text{CI}}/\Sigma$ (metabolic), $[\text{AAb}]$ (neurovascular) \\
\hline
\end{tabular}
\end{table}

The formal analysis both confirms and refines the qualitative classification. The major refinement is the identification of \emph{subtype dependence}: a mechanism's role in the causal hierarchy depends on which disease attractor the patient occupies. Complex~I reduction is a trigger and a lock for the metabolic subtype but merely an amplifier for the immune-dominant subtype. This context-dependent classification is invisible to qualitative reasoning but emerges naturally from the mathematical structure of the coupled ODE system.

\subsection*{Limitations and Next Steps}

The analyses in this chapter rest on several assumptions that must be tested:

\begin{enumerate}
    \item \textbf{Reduced model validity}: Trigger thresholds and lock removal outcomes were derived from reduced subsystem models, not the full 67-variable system. Emergent behaviors arising from the coupling of all subsystems may shift thresholds or change qualitative outcomes. Numerical simulation of the full system with realistic parameter ranges is the essential next step.

    \item \textbf{Parameter uncertainty}: Many model parameters (Tables~\ref{tab:aab-parameters}--\ref{tab:safe-mode-parameters}) are estimated rather than measured. The qualitative predictions (which mechanisms are triggers, locks, or amplifiers) are expected to be robust to moderate parameter uncertainty, but the quantitative thresholds may shift substantially. Bayesian uncertainty propagation through the full model would quantify this robustness.

    \item \textbf{Discrete patient subtyping}: The model predicts discrete attractors, but real patients may occupy intermediate positions (near separatrices, transitioning between attractors, or in mixed states). The sharpness of the subtype boundaries in practice is an empirical question addressable through multi-omics clustering studies.

    \item \textbf{Treatment as parameter restoration}: The priority score (Equation~\ref{eq:priority-score}) assumes that treatments act by restoring disease parameters toward healthy values. In practice, treatments have off-target effects, pharmacokinetic variability, and dose--response nonlinearities not captured in the current analysis. Integration with the pharmacokinetic--pharmacodynamic framework of Chapter~\ref{ch:predictive-applications} (Section~\ref{sec:dosing-optimization}) would improve clinical relevance.

    \item \textbf{Temporal sequencing of interventions}: The minimum intervention set analysis (Section~\ref{sec:minimum-intervention-sets}) identifies \emph{which} parameters must be restored simultaneously but does not address \emph{temporal ordering}. In practice, some interventions must precede others (e.g., reducing inflammation before attempting epigenetic reversal). Optimal temporal sequencing requires dynamic optimization over the ODE trajectories, an extension left for future work.
\end{enumerate}

The critical next step is numerical validation of the full 67-variable system, which will either confirm the qualitative predictions derived here from reduced models or reveal emergent behaviors not captured by subsystem analysis. Chapter~\ref{ch:predictive-applications} translates these findings into clinical decision support tools; the formal hierarchy analysis provides the mechanistic foundation that those tools require.
